%% toc: P15 159 159.7,159.10
\section{A probability space without a regular conditional probability}
\label{sec:annex-p15-w02c-p15-008-no-rcp}

Throughout, \(\lambda\) is Lebesgue measure on \(\mathcal B(\mathbb R)\)
from \cref{thm:p15-lebesgue-unit-interval},
\((\Omega_U,\mathcal F_U,\mathbb U)\) is the uniform probability space
there, and \(L,L_0,L_1\) are the sets of
\cref{def:p15-w02c-p15-008-groups}. In the last two items complements of
subsets of \(\Omega_U=(0,1)\) are taken in \(\Omega_U\).

\begin{annexitem}{annex:p15-w02c-p15-008-groups}{thm:p15-w02c-p15-008-groups}{Parity of the integer part and small odd elements}
\begin{proof}
\paragraph{Irrationality of \(\sqrt2\).}
Suppose \(\sqrt2=p/q\) with positive integers \(p,q\), and take \(q\)
least among all such representations. Then \(p^2=2q^2\), so \(p^2\) is
even; since the square of an odd integer is odd, \(p\) is even, say
\(p=2r\). Then \(4r^2=2q^2\), so \(q^2=2r^2\) and \(q\) is even, say
\(q=2s\). Now \(\sqrt2=r/s\) with \(0<s<q\), contradicting the choice
of \(q\). Hence \(\sqrt2\) is irrational.

\paragraph{Part (i).}
For integers \(m,n,m',n'\),
\[
 \begin{gathered}
 (m+n\sqrt2)+(m'+n'\sqrt2)=(m+m')+(n+n')\sqrt2,\\
 -(m+n\sqrt2)=(-m)+(-n)\sqrt2,
 \end{gathered}
\]
and \(0=0+0\sqrt2\). So \(L\) contains \(0\) and is closed under
addition and negation, hence is an additive subgroup. The same
computations with \(m,m'\) replaced by \(2m,2m'\) show that \(L_0\) is a
subgroup, since \(2m+2m'\) and \(-2m\) are even. The identity
\(2m+1+n\sqrt2=(2m+n\sqrt2)+1\) gives \(L_1=L_0+1\). Every integer is
even or odd, so \(L=L_0\cup L_1\).

The union is disjoint. If \(2m+n\sqrt2=2m'+1+n'\sqrt2\), then
\((n-n')\sqrt2=2(m'-m)+1\). If \(n=n'\), the odd integer \(2(m'-m)+1\)
would be \(0\), which is impossible. If \(n\neq n'\), then
\(\sqrt2=(2(m'-m)+1)/(n-n')\) would be rational, which is impossible by
the first paragraph.

For \(c=2m+1+n\sqrt2\in L_1\) and \(b=2m'+n'\sqrt2\in L_0\),
\(c+b=2(m+m')+1+(n+n')\sqrt2\in L_1\). Conversely each \(c\in L_1\) equals
\(c+0\) with \(0\in L_0\). Hence \(L_1+L_0=L_1\).

\paragraph{Small positive elements of \(L_0\).}
The set \(L_0\cap(0,\infty)\) contains \(2\); let \(d\geq0\) be its
infimum. Suppose \(d>0\).

First, \(d\in L_0\). Otherwise \(d\) is not attained, so the definition
of the infimum gives \(b_2\in L_0\) with \(d<b_2<2d\), and then
\(b_1\in L_0\) with \(d<b_1<b_2\). The difference \(b_2-b_1\) lies in
\(L_0\), because \(L_0\) is a subgroup, and satisfies
\(0<b_2-b_1<2d-d=d\). This contradicts the definition of \(d\).

Next, \(L_0\subseteq d\mathbb Z\). For \(g\in L_0\) let
\(k=\lfloor g/d\rfloor\). Then \(g-kd\in L_0\) and \(0\leq g-kd<d\).
No member of \(L_0\) lies in \((0,d)\), so \(g-kd=0\).

In particular \(2=kd\) and \(\sqrt2=ld\) for some integers \(k,l\), and
\(k\neq0\). Then \(\sqrt2=2l/k\) would be rational, a contradiction.
Hence \(d=0\): for every \(\varepsilon>0\) there is \(b\in L_0\) with
\(0<b<\varepsilon\).

\paragraph{Part (ii).}
Let \(\delta>0\) and choose \(b\in L_0\) with \(0<b<\delta\). The set of
integers \(k\) with \(kb>-1-\delta\) is nonempty and bounded below, by the
Archimedean property; let \(k\) be its least member. Then
\(-1-\delta<kb\leq-1-\delta+b<-1\), so
\[
 -\delta<1+kb<0<\delta .
\]
Since \(kb\in L_0\) and \(L_1=L_0+1\), the number \(1+kb\) lies in
\(L_1\cap(-\delta,\delta)\).
\end{proof}
\end{annexitem}

\begin{annexitem}{annex:p15-w02c-p15-008-steinhaus}{thm:p15-w02c-p15-008-steinhaus}{An interval in which the set fills more than three quarters}
\begin{proof}
\paragraph{The cover description of \(\lambda\).}
By \cref{thm:p15-lebesgue-unit-interval}, \(\lambda\) is the measure of
\cref{thm:p15-stieltjes-construction} for the identity function \(x\mapsto x\). That theorem
obtains it from the premeasure \(\overline m\) on the interval algebra
\(\mathcal A(\mathcal S_{\mathbb R})\) given by
\cref{thm:p15-semialgebra-extension}, and
\cref{thm:p15-premeasure-extension} identifies the resulting measure as
the restriction to \(\sigma(\mathcal A(\mathcal S_{\mathbb R}))
=\mathcal B(\mathbb R)\) of the outer measure \(\lambda^*\) of
\cref{def:p15-cover-outer-measure}. Hence, for Borel \(F\),
\[
 \lambda(F)=\inf\Bigl\{\sum_{n\geq1}\overline m(A_n):
           A_n\in\mathcal A(\mathcal S_{\mathbb R}),\quad
           F\subseteq\bigcup_{n\geq1}A_n\Bigr\}.
\]

\paragraph{A good interval.}
Let \(0<\lambda(F)<\infty\). Then \(\tfrac43\lambda(F)>\lambda(F)\), so
there are \(A_n\in\mathcal A(\mathcal S_{\mathbb R})\) with
\(F\subseteq\bigcup_nA_n\) and
\(\sum_n\overline m(A_n)<\tfrac43\lambda(F)\). Each \(A_n\) is a finite
disjoint union of members of \(\mathcal S_{\mathbb R}\), and
\(\overline m(A_n)\) is the sum of their \(m\)-values
(\cref{thm:p15-semialgebra-extension}). List the nonempty members of all
these finite unions as one sequence \(J_1,J_2,\ldots\) (finite or
infinite). Then \(F\subseteq\bigcup_kJ_k\) and
\[
 \sum_k m(J_k)<\tfrac43\lambda(F)<\infty .
\]
Every \(m(J_k)\) is finite. A nonempty member of
\(\mathcal S_{\mathbb R}\) with an infinite endpoint has \(m\)-value
\(\infty\) for the identity function, so \(J_k=(a_k,b_k]\) with real
\(a_k<b_k\), and \(\lambda(J_k)=b_k-a_k=m(J_k)\).

By monotonicity and countable subadditivity
(\cref{thm:p15-measure-laws}),
\(\lambda(F)\leq\sum_k\lambda(F\cap J_k)\). If
\(\lambda(J_k)\geq\tfrac43\lambda(F\cap J_k)\) held for every \(k\),
summing would give
\(\sum_km(J_k)\geq\tfrac43\sum_k\lambda(F\cap J_k)\geq\tfrac43\lambda(F)\),
contrary to the choice of the cover. So there is a \(k\) with
\[
 \lambda(F\cap J_k)>\tfrac34\lambda(J_k).
\]
Write \(J=(a,b]\) for this \(J_k\), \(\ell=b-a>0\), and
\(F_J=F\cap J\). Then \(F_J\) is Borel and \(\lambda(F_J)>\tfrac34\ell\).

\paragraph{Differences.}
Put \(\delta=\ell/2\) and let \(|t|<\delta\). The intervals \(J\) and
\(J+t=(a+t,b+t]\) overlap, because \(|t|<\ell\), so their union is the
interval \((\min(a,a+t),\max(b,b+t)]\), of length \(\ell+|t|\) by
\cref{thm:p15-lebesgue-unit-interval}. The set \(F_J+t\) is Borel with
\(\lambda(F_J+t)=\lambda(F_J)\), by translation invariance in the same theorem,
and \(F_J\cup(F_J+t)\subseteq J\cup(J+t)\).
Suppose \(F_J\cap(F_J+t)=\varnothing\). Then additivity and monotonicity give
\[
 2\lambda(F_J)=\lambda\bigl(F_J\cup(F_J+t)\bigr)\leq\ell+|t|<\tfrac32\ell
           <2\lambda(F_J),
\]
a contradiction. So there is \(x\in F_J\cap(F_J+t)\), that is, \(x\in F_J\)
and \(x=y+t\) with \(y\in F_J\). Then \(t=x-y\in F_J-F_J\subseteq F-F\). As
\(t\) was arbitrary with \(|t|<\delta\), \((-\delta,\delta)\subseteq F-F\).
\end{proof}
\end{annexitem}

\begin{annexitem}{annex:p15-w02c-p15-008-thick-set}{thm:p15-w02c-p15-008-thick-set}{Unique coset representation and a forbidden difference}
\begin{proof}
\paragraph{The selector.}
Since \(L\) is a subgroup, \(x\sim y\) if and only if \(x-y\in L\) is an
equivalence relation on \(\mathbb R\), and its classes are the cosets
\(x+L\). These cosets are nonempty and pairwise disjoint. The axiom of
choice gives a set \(V\) containing exactly one point of each of them.

\paragraph{Unique representation.}
Every \(x\in\mathbb R\) can be written as \(x=v+l\) with \(v\in V\) and
\(l\in L\), in exactly one way. For existence, let \(v\) be the point of
\(V\) in \(x+L\); then \(v=x+l'\) with \(l'\in L\), and \(x=v+(-l')\).
For uniqueness, if \(v+l=v'+l'\) with \(v,v'\in V\) and \(l,l'\in L\),
then \(v-v'=l'-l\in L\), so \(v\) and \(v'\) lie in the same coset; as
\(V\) has one point in each coset, \(v=v'\), and then \(l=l'\).

\paragraph{The complement.}
Let \(x=v+l\) be the representation. If \(l\in L_0\), then \(x\in H\).
If \(x\in H\), say \(x=v'+b\) with \(v'\in V\) and \(b\in L_0\subseteq L\),
uniqueness gives \(l=b\in L_0\). So \(x\in H\) exactly when \(l\in L_0\).
By part~(i) of \cref{thm:p15-w02c-p15-008-groups}, \(L\) is the disjoint
union of \(L_0\) and \(L_1\), so \(x\notin H\) exactly when \(l\in L_1\).
Hence \(\mathbb R\setminus H=V+L_1=V+(L_0+1)=H+1\).

\paragraph{A forbidden difference.}
We show \((H-H)\cap L_1=\varnothing\). Suppose
\((v+b)-(v'+b')=c\in L_1\) with \(v,v'\in V\) and \(b,b'\in L_0\). Then
\(v-v'=c+(b'-b)\in L_1+L_0=L_1\subseteq L\), so \(v=v'\) as before, and
\(c=b-b'\in L_0\). This contradicts \(L_0\cap L_1=\varnothing\).

\paragraph{Borel subsets of \(H\).}
Let \(F\subseteq H\) be Borel and suppose \(\lambda(F)>0\). The sets
\(F\cap(-k,k]\), \(k\geq1\), increase to \(F\), so continuity from below
(\cref{thm:p15-measure-laws}) gives a \(k\) with
\(\lambda(F\cap(-k,k])>0\); this measure is at most \(2k\).
\Cref{thm:p15-w02c-p15-008-steinhaus} gives \(\delta>0\) with
\[
 (-\delta,\delta)\subseteq\bigl(F\cap(-k,k]\bigr)-\bigl(F\cap(-k,k]\bigr)
                \subseteq H-H .
\]
Part~(ii) of \cref{thm:p15-w02c-p15-008-groups} gives
\(c\in L_1\cap(-\delta,\delta)\), so \(c\in(H-H)\cap L_1\), contrary to
the preceding paragraph. Hence \(\lambda(F)=0\).

\paragraph{Borel subsets of the complement.}
Let \(F\subseteq\mathbb R\setminus H=H+1\) be Borel. Then \(F-1\) is
Borel and \(F-1\subseteq H\), so \(\lambda(F-1)=0\). Translation
invariance (\cref{thm:p15-lebesgue-unit-interval}) applied to \(F-1\) and
\(t=1\) gives \(\lambda(F)=\lambda(F-1)=0\).

\paragraph{The trace on \((0,1)\).}
Every member of \(\mathcal F_U\) is a Borel subset of \(\mathbb R\)
contained in \((0,1)\), as shown in the proof of
\cref{thm:p15-lebesgue-unit-interval}, and \(\mathbb U\) is \(\lambda\)
on it. A member contained in \(H_U\subseteq H\) or in
\(H_U^c\subseteq\mathbb R\setminus H\) therefore has \(\mathbb U\)-measure
zero. If \(H_U\) belonged to \(\mathcal F_U\), so would
\(H_U^c=(0,1)\setminus H_U\); applying this to \(H_U\subseteq H_U\) and
\(H_U^c\subseteq H_U^c\) would give
\(1=\mathbb U((0,1))=\mathbb U(H_U)+\mathbb U(H_U^c)=0\). Hence
\(H_U\notin\mathcal F_U\).
\end{proof}
\end{annexitem}

\begin{annexitem}{annex:p15-w02c-p15-008-model}{thm:p15-w02c-p15-008-model}{Representations, null symmetric differences and the two versions}
\begin{proof}
\paragraph{Part (i).}
For \(B\in\mathcal F_U\), \(B=(B\cap H_U)\cup(B\cap H_U^c)\in\mathcal F_H\).
Also \(H_U=(\Omega_U\cap H_U)\cup(\varnothing\cap H_U^c)\in\mathcal F_H\),
and \(\Omega_U\in\mathcal F_H\). For
\(D=(B_1\cap H_U)\cup(B_2\cap H_U^c)\),
\[
 \Omega_U\setminus D=\bigl((\Omega_U\setminus B_1)\cap H_U\bigr)
                     \cup\bigl((\Omega_U\setminus B_2)\cap H_U^c\bigr),
\]
because a point of \(H_U\) lies outside \(D\) exactly when it lies
outside \(B_1\), and a point of \(H_U^c\) exactly when it lies outside
\(B_2\). For a sequence \(D_n=(B_{1n}\cap H_U)\cup(B_{2n}\cap H_U^c)\),
\[
 \bigcup_nD_n=\Bigl(\bigl(\textstyle\bigcup_nB_{1n}\bigr)\cap H_U\Bigr)
              \cup\Bigl(\bigl(\textstyle\bigcup_nB_{2n}\bigr)\cap H_U^c\Bigr).
\]
\(\mathcal F_U\) is closed under complements in \(\Omega_U\) and countable
unions, so \(\mathcal F_H\) is a sigma-algebra.

\paragraph{A null-difference fact.}
Let \(B,B'\in\mathcal F_U\). If \(B\cap H_U=B'\cap H_U\), then
\(\mathbb U(B)=\mathbb U(B')\); the same holds with \(H_U^c\) in place of
\(H_U\). Indeed, the symmetric difference \(B\triangle B'\) belongs to
\(\mathcal F_U\) and does not meet \(H_U\), so it is contained in
\(H_U^c\) and has measure zero by
\cref{thm:p15-w02c-p15-008-thick-set}. Then
\(\mathbb U(B)=\mathbb U(B\cap B')+\mathbb U(B\setminus B')
=\mathbb U(B\cap B')\), and likewise for \(B'\). The case of \(H_U^c\) is
the same with the roles of \(H_U\) and \(H_U^c\) exchanged.

\paragraph{Part (ii).}
If \((B_1\cap H_U)\cup(B_2\cap H_U^c)=(B_1'\cap H_U)\cup(B_2'\cap H_U^c)\),
intersecting both sides with \(H_U\) gives \(B_1\cap H_U=B_1'\cap H_U\),
and with \(H_U^c\) gives \(B_2\cap H_U^c=B_2'\cap H_U^c\). The
null-difference fact gives \(\mathbb U(B_1)=\mathbb U(B_1')\) and
\(\mathbb U(B_2)=\mathbb U(B_2')\), so \(P_H(D)\) does not depend on the
representation. Its values lie in \([0,1]\); \(P_H(\varnothing)=0\) and
\(P_H(\Omega_U)=1\).

Let \(D_n=(B_{1n}\cap H_U)\cup(B_{2n}\cap H_U^c)\) be pairwise disjoint.
Put \(B_{1n}^*=B_{1n}\setminus\bigcup_{m<n}B_{1m}\). These sets lie in
\(\mathcal F_U\), are pairwise disjoint, and have the same union as the
\(B_{1n}\). For \(m<n\),
\(B_{1n}\cap B_{1m}\cap H_U\subseteq D_n\cap D_m=\varnothing\), so
\(B_{1n}^*\cap H_U=B_{1n}\cap H_U\) and, by the null-difference fact,
\(\mathbb U(B_{1n}^*)=\mathbb U(B_{1n})\). Define \(B_{2n}^*\) in the same
way; the same argument with \(H_U^c\) gives
\(\mathbb U(B_{2n}^*)=\mathbb U(B_{2n})\). Using the union formula of
part~(i) and countable additivity of \(\mathbb U\),
\[
 \begin{aligned}
 P_H\Bigl(\bigcup_nD_n\Bigr)
  &=\tfrac12\Bigl(\sum_n\mathbb U(B_{1n}^*)+\sum_n\mathbb U(B_{2n}^*)\Bigr)\\
  &=\sum_n\tfrac12\bigl(\mathbb U(B_{1n})+\mathbb U(B_{2n})\bigr)
  =\sum_nP_H(D_n),
 \end{aligned}
\]
where the series have nonnegative terms. So \(P_H\) is a probability
measure. For \(B\in\mathcal F_U\), the representations
\((B\cap H_U)\cup(B\cap H_U^c)\), \((B\cap H_U)\cup(\varnothing\cap H_U^c)\)
and \((\varnothing\cap H_U)\cup(B\cap H_U^c)\) give
\(P_H(B)=\mathbb U(B)\) and
\(P_H(B\cap H_U)=P_H(B\cap H_U^c)=\tfrac12\mathbb U(B)\).

\paragraph{Part (iii).}
Let \(B\in\mathcal F_U\). The indicator \(\mathbf1_B\) is
\(\mathcal F_U\)-measurable and bounded, and for \(A\in\mathcal F_U\) the
two sides of the defining identity of
\cref{def:p15-general-conditional-expectation} for
\(E(\mathbf1_B\mid\mathcal F_U)\) are the same integral
\(\int_A\mathbf1_B\,dP_H\). So \(\mathbf1_B\) is a version of
\(P_H(B\mid\mathcal F_U)\).

The constant \(\tfrac12\) is \(\mathcal F_U\)-measurable. For
\(A\in\mathcal F_U\), by part~(ii),
\[
 \int_A\tfrac12\,dP_H=\tfrac12P_H(A)=\tfrac12\mathbb U(A)
   =P_H(A\cap H_U)=\int_A\mathbf1_{H_U}\,dP_H .
\]
So \(\tfrac12\) is a version of \(P_H(H_U\mid\mathcal F_U)\). If \(Z\) is
any version, \cref{thm:p15-conditional-existence} gives \(Z=\tfrac12\)
almost everywhere; the set \(\{Z\neq\tfrac12\}\) belongs to
\(\mathcal F_U\) because \(Z\) is \(\mathcal F_U\)-measurable, and it is
\(P_H\)-null.
\end{proof}
\end{annexitem}

\begin{annexitem}{annex:p15-w02c-p15-008-no-rcp}{thm:p15-w02c-p15-008-no-rcp}{Forced point masses}
\begin{proof}
Let \(\mathbb Q_0=\mathbb Q\cap(0,1)\) and
\(\mathcal C=\{(0,q]:q\in\mathbb Q_0\}\cup\{H_U\}\), a countable subfamily
of \(\mathcal F_H\). For \(C\in\mathcal C\) let \(h_C\) be the function
\(\mathbf1_{(0,q]}\) when \(C=(0,q]\), and the constant \(\tfrac12\) when
\(C=H_U\).

\paragraph{Key step.}
No function \(\mu:\Omega_U\times\mathcal F_H\to[0,1]\) has both of the
following properties:
\begin{enumerate}
\item[(a)] for each \(C\in\mathcal C\) there is a \(P_H\)-null set
      \(N_C\in\mathcal F_H\) with \(\mu(\omega,C)=h_C(\omega)\) for
      \(\omega\notin N_C\);
\item[(b)] there is a \(P_H\)-null set \(N\in\mathcal F_H\) such that
      \(\mu(\omega,\cdot)\) is a probability measure on
      \(\mathcal F_H\) for \(\omega\notin N\).
\end{enumerate}
Suppose \(\mu\) had both. The set \(M=N\cup\bigcup_{C\in\mathcal C}N_C\)
is a countable union of null sets, so \(P_H(M)=0\) by countable
subadditivity (\cref{thm:p15-measure-laws}). As \(P_H(\Omega_U)=1\),
there is \(\omega\in\Omega_U\setminus M\). Fix it and put
\(\nu=\mu(\omega,\cdot)\), a probability measure on \(\mathcal F_H\) with
\[
 \nu((0,q])=\mathbf1_{(0,q]}(\omega)\quad(q\in\mathbb Q_0),\qquad
 \nu(H_U)=\tfrac12 .
\]
Since \(0<\omega<1\), density of the rationals gives, inductively, rationals
\(q_n\) with \(\omega<q_n<\min(q_{n-1},\omega+1/n)\) (with \(q_0=1\)) and
rationals \(r_n\) with \(\max(r_{n-1},\omega-1/n)<r_n<\omega\) (with
\(r_0=0\)). All lie in \(\mathbb Q_0\). Then \(\nu((0,q_n])=1\) and
\(\nu((0,r_n])=0\) for every \(n\). The sets \((0,q_n]\) decrease to
\((0,\omega]\), and the sets \((0,r_n]\) increase to \((0,\omega)\); all
belong to \(\mathcal F_U\subseteq\mathcal F_H\). Continuity from above
and from below (\cref{thm:p15-measure-laws}) give
\(\nu((0,\omega])=1\) and \(\nu((0,\omega))=0\), hence
\[
 \nu(\{\omega\})=\nu((0,\omega])-\nu((0,\omega))=1 .
\]
If \(\omega\in H_U\), then \(\{\omega\}\subseteq H_U\) and
\(\nu(H_U)\geq1\). If \(\omega\notin H_U\), then
\(H_U\subseteq\Omega_U\setminus\{\omega\}\) and
\(\nu(H_U)\leq1-\nu(\{\omega\})=0\). Either way \(\nu(H_U)\neq\tfrac12\), a
contradiction. This proves the key step.

\paragraph{Part (i).}
Let \(\mu\) satisfy \cref{def:p15-w02b-p15-008-rcd} for the identity
\(X\), read with the exceptional set of
\cref{def:p15-w02b-p15-008-exceptional-set}. Property~(b) is condition~(ii)
of that definition. For \(C\in\mathcal C\), \(\{X\in C\}=C\), so
\(\mu(\cdot,C)\) is a version of \(P_H(C\mid\mathcal F_U)\). By
part~(iii) of \cref{thm:p15-w02c-p15-008-model}, \(h_C\) is also a
version, and \cref{thm:p15-conditional-existence} gives a \(P_H\)-null
set outside which the two agree; the set
\(\{\mu(\cdot,C)\neq h_C\}\) belongs to \(\mathcal F_U\subseteq\mathcal F_H\)
since both functions are \(\mathcal F_U\)-measurable. This is
property~(a). The key step excludes such a \(\mu\).

\paragraph{Part (ii).}
Let \(K\) be a probability kernel from \((\Omega_U,\mathcal F_U)\) to
\((\Omega_U,\mathcal F_H)\) with \(K(\cdot,A)=P_H(X\in A\mid\mathcal F_U)\)
almost everywhere for each \(A\in\mathcal F_H\). Property~(b) holds with
\(N=\varnothing\). For \(C\in\mathcal C\), \(K(\cdot,C)\) agrees outside a
null set with a version of \(P_H(C\mid\mathcal F_U)\), and that version
agrees with \(h_C\) outside a null set by
\cref{thm:p15-conditional-existence} and part~(iii) of
\cref{thm:p15-w02c-p15-008-model}. Both \(K(\cdot,C)\) and \(h_C\) are
\(\mathcal F_U\)-measurable, so \(\{K(\cdot,C)\neq h_C\}\) is an
\(\mathcal F_U\)-measurable set contained in the union of the two null
sets, hence null. This is property~(a), and the key step excludes \(K\).

\paragraph{Part (iii).}
Suppose \(\phi:\Omega_U\to E\), with \(E\) a Borel subset of
\(\mathbb R\), were a measurable isomorphism for \(\mathcal F_H\) as in
\cref{thm:p15-standard-conditional-law}. The identity
\(X:(\Omega_U,\mathcal F_H,P_H)\to(\Omega_U,\mathcal F_H)\) is measurable,
so that theorem would give a regular conditional law of \(X\) given
\(\mathcal F_U\) in the sense of \cref{def:p15-regular-conditional-law},
that is, a probability kernel excluded by part~(ii).

\paragraph{Part (iv).}
The identity \(Y:(\Omega_U,\mathcal F_H)\to(\Omega_U,\mathcal F_U)\) is
measurable, since \(Y^{-1}(B)=B\in\mathcal F_U\subseteq\mathcal F_H\). For
\(B\in\mathcal F_U\), \(\omega\mapsto\delta(\omega,B)=\mathbf1_B(\omega)\)
is \(\mathcal F_U\)-measurable. For fixed \(\omega\),
\(B\mapsto\mathbf1_B(\omega)\) has value \(1\) at \(\Omega_U\), and is
countably additive: \(\omega\) lies in at most one member of a disjoint
sequence, and in their union exactly when it lies in one of them. So
\(\delta\) is a probability kernel from \((\Omega_U,\mathcal F_U)\) to
\((\Omega_U,\mathcal F_U)\). By part~(iii) of
\cref{thm:p15-w02c-p15-008-model}, \(\delta(\cdot,B)=\mathbf1_B\) is a
version of \(P_H(Y\in B\mid\mathcal F_U)=P_H(B\mid\mathcal F_U)\), so the
identity of \cref{def:p15-regular-conditional-law} holds at every
outcome. Finally, \cref{thm:p15-real-conditional-law} applies to every
real random variable on \((\Omega_U,\mathcal F_H,P_H)\) and every
sub-sigma-algebra, in particular \(\mathcal F_U\).
\end{proof}
\end{annexitem}
